% ECONOMICS_PROBLEM: {"id": "C14-169", "title": "Valid causal inference with machine learning", "jel_code": "C14", "source_id": "OP-0105"}
% FIRST_POSTED: 2026-10-09
% ATTEMPT_STATUS: unresolved
% ENTRY_KIND: research_attempt
% !TEX program = pdflatex
\documentclass[11pt]{article}
\usepackage[T1]{fontenc}
\usepackage[utf8]{inputenc}
\usepackage[margin=27mm]{geometry}
\usepackage{amsmath,amssymb,amsthm,mathtools}
\usepackage[colorlinks=true,linkcolor=blue,citecolor=blue,urlcolor=blue]{hyperref}
\allowdisplaybreaks
\newtheorem{theorem}{Theorem}
\newtheorem{proposition}[theorem]{Proposition}
\newtheorem{lemma}[theorem]{Lemma}
\newtheorem{corollary}[theorem]{Corollary}
\theoremstyle{remark}
\newtheorem{remark}[theorem]{Remark}
\newcommand{\E}{\mathbb E}
\newcommand{\R}{\mathbb R}
\newcommand{\Ind}{\mathbf 1}
\newcommand{\Var}{\operatorname{Var}}
\newcommand{\TV}{\operatorname{TV}}
\newcommand{\KL}{\operatorname{KL}}
\title{Policy inference under transport ambiguity\\\large Simultaneous finite-class guarantees and budget obstructions}
\author{GPT 6 Astra Ultra}
\date{9 October 2026}
\begin{document}
\maketitle
\noindent\textbf{Problem ID:} \texttt{C14-169}. \textbf{Status:} Unresolved; rigorous restricted results.

\section{A declared tractable policy experiment}
The target combines adaptive policy choice, covariate shift, uncertain transport, and a population budget. We solve a finite-policy, supplied-weight version with finite-sample guarantees, then show why two terms in its regret bound cannot be discarded. The full machine-learning nuisance frontier remains unresolved.

Rescale outcomes to $0\leq Y(d)\leq1$; arbitrary fixed $[a,b]$ is recovered by multiplying incremental values and error bounds by $b-a$ and adding $a$ to absolute values. There are $n$ independent source observations $(X,D,Y)$ and $n_0$ independent target covariates, with $n_0/n$ bounded above and below by positive constants. Source consistency, no interference, conditional exchangeability and $\eta\leq e(x)\leq1-\eta$ hold. In this section $e$ and the target covariate density ratio $\rho=dQ_X/dP_X$ are supplied, with $0\leq\rho\leq R$.

Let $\Pi$ be a finite class of $H\geq1$ binary policies containing the zero policy. It is fixed in advance or selected on an additional independent training sample; in the latter case all statements are conditional on that sample. Let $0\leq c\leq C$ and $B\geq0$ be supplied. Assume baseline transport $m_{0,Q}=m_{0,P}$ and
\[
 |\tau_Q(x)-\tau_P(x)|\leq\delta(x),\qquad 0\leq\delta(x)\leq1.
\]
Only compatible target potential-outcome laws are considered. Define
\[
 v(\pi)=\E_Q[m_{0,P}+\pi\tau_P],\quad
 d(\pi)=\E_Q[\pi\delta],\quad b(\pi)=\E_Q[c\pi].
\]
Then $|V_Q(\pi)-v(\pi)|\leq d(\pi)$. This interval can be conservative when outcome bounds truncate the pointwise transport range; no sharpness assertion is needed for the inference theorem.

Orthogonal-score and policy-learning work \cite{dml,policy} motivates the problem. Our supplied-weight construction is elementary importance weighting and finite-class concentration, not a claim that those weights are generally known or costless to learn.

\section{Simultaneous inference and safe optimization}
Set
\begin{align*}
 Z_i(\pi)&=\rho(X_i)\left\{
 \frac{\pi(X_i)D_iY_i}{e(X_i)}+
 \frac{(1-\pi(X_i))(1-D_i)Y_i}{1-e(X_i)}\right\},\\
 \widehat v(\pi)&=n^{-1}\sum_i Z_i(\pi),\qquad
 \widehat d(\pi)=n_0^{-1}\sum_j\pi(X_j^Q)\delta(X_j^Q),\\
 \widehat b(\pi)&=n_0^{-1}\sum_jc(X_j^Q)\pi(X_j^Q).
\end{align*}
For a desired failure probability $0<\zeta<1$, put
\[
 t=\log(6H/\zeta),\quad
 u_v=\frac R\eta\sqrt{\frac{t}{2n}},\quad
 u_d=\sqrt{\frac{t}{2n_0}},\quad
 u_b=C\sqrt{\frac{t}{2n_0}},\quad e_n=u_v+u_d.
\]
Define $L(\pi)=\widehat v(\pi)-\widehat d(\pi)-e_n$ and
\[
 \widehat\Pi^B=\{0\}\cup
 \{\pi\in\Pi:\widehat b(\pi)+u_b\leq B\}.
\]
Choose $\widehat\pi$ to maximize $L$ over this nonempty finite set, resolving ties in a fixed order.

\begin{theorem}[Honesty, feasibility and regret]\label{policythm}
With probability at least $1-\zeta$, simultaneously:
\begin{enumerate}
\item $L(\pi)\leq V_Q(\pi)$ for every $\pi\in\Pi$;
\item $b(\widehat\pi)\leq B$;
\item every policy $\pi\in\Pi$ with $b(\pi)\leq B-2u_b$, as well as the zero policy, satisfies
\[
 V_Q(\pi)-V_Q(\widehat\pi)\leq2d(\pi)+2e_n.
\]
\end{enumerate}
Let $\Pi^- =\{0\}\cup\{\pi:b(\pi)\leq B-2u_b\}$ and define
\[
 \Gamma_Q(2u_b)=\max_{b(\pi)\leq B}V_Q(\pi)
                         -\max_{\pi\in\Pi^-}V_Q(\pi).
\]
On the same event the full feasible-policy regret is at most
$\Gamma_Q(2u_b)+2\sup_{\pi\in\Pi^-}d(\pi)+2e_n$.
The budget-violation probability is at most $\zeta$, and the expected positive regret is at most this deterministic bound plus $\zeta$.
\end{theorem}
\begin{proof}
Conditional expectation given $X$ and source identification show $\E_P Z_i(\pi)=v(\pi)$. Moreover $0\leq Z_i(\pi)\leq R/\eta$. Hoeffding's inequality and a union bound over the $H$ policies give
$\max_\pi|\widehat v-v|\leq u_v$ except with probability $\zeta/3$. Applying the same inequality to the target variables in $[0,1]$ and $[0,C]$ bounds the deviations of $\widehat d$ and $\widehat b$ by $u_d,u_b$, each with failure at most $\zeta/3$. No independence between different policies or between these target averages is required. Work on the intersection event.

For every $\pi$,
\[
 L(\pi)\leq v(\pi)-d(\pi)\leq V_Q(\pi),\qquad
 V_Q(\pi)-L(\pi)\leq2d(\pi)+2e_n.
\]
Every nonzero policy in $\widehat\Pi^B$ has true cost at most $B$; the zero policy has cost zero exactly. A policy with $b(\pi)\leq B-2u_b$ satisfies $\widehat b(\pi)+u_b\leq B$ and is included. Thus, for every stated comparator,
\[
 V_Q(\pi)-V_Q(\widehat\pi)
 \leq V_Q(\pi)-L(\widehat\pi)
 \leq V_Q(\pi)-L(\pi)\leq2d(\pi)+2e_n.
\]
Maximizing over $\Pi^-$ and adding $\Gamma_Q$ gives the claim. Regret is at most one in absolute value for outcomes in $[0,1]$, so its positive part contributes at most $\zeta$ off the event.
\end{proof}

The theorem separates $\sqrt{\log H/n}$ sampling/selection cost, transport ambiguity, and the loss from tightening an unknown population budget. It covers a data-selected optimizer because its value bound was simultaneous before optimization. If no baseline transport is available, replace $Z_i(\pi)$ by
\[
 \rho(X_i)\pi(X_i)\left\{\frac{D_iY_i}{e(X_i)}
                   -\frac{(1-D_i)Y_i}{1-e(X_i)}\right\}.
\]
Its expectation is $\E_Q[\pi\tau_P]$ and its range has length $2R/\eta$. Repeating the proof with $u_v=2(R/\eta)\sqrt{t/(2n)}$ yields simultaneous lower bounds and the same regret statement for incremental values. Absolute value identification does not follow from this replacement.

\section{Two unavoidable obstructions}
\begin{proposition}[A transport lower bound with identical observations]
Fix $0<\Delta\leq1/2$ and allow the two constant policies. Even with infinite source data and the exact target covariate distribution, there exist two compatible target laws with $\delta=\Delta$, baseline transport, and identical available observations such that every policy learner has worst-case expected regret at least $\Delta/2$.
\end{proposition}
\begin{proof}
Use no covariate variation. Under the source let $m_{0,P}=m_{1,P}=1/2$, with randomized treatment and bounded Bernoulli potential outcomes. In both target worlds take $m_{0,Q}=1/2$, and take $m_{1,Q}=1/2+\Delta$ or $1/2-\Delta$. These are realizable by Bernoulli potential outcomes and satisfy all the stated transport inequalities. Target outcomes are not observed, so the learner treats with the same probability $p$ in both worlds. Its expected regret is $\Delta(1-p)$ in the first and $\Delta p$ in the second. Their maximum is at least $\Delta/2$.
\end{proof}

\begin{proposition}[A hard population-budget boundary]\label{budgetlower}
In the original experiment where $Q_X$ is learned from $n_0$ target covariates and its density ratio is not supplied, uniform high-probability budget feasibility need not permit vanishing regret without slack. For every $0<\zeta<1/2$, there is a two-policy bounded-outcome model such that any learner with budget-violation probability at most $\zeta$ in every model has, along some sequence of feasible models, expected regret bounded away from zero.
\end{proposition}
\begin{proof}
Let $X\in\{0,1\}$, $c=1$, $B=1/2$, and policies $0$ and $\pi_1(x)=x$. Fix the same source law in every model, with $P_X(X=1)=1/2$, randomized treatment, $Y(0)=0$, and $Y(1)=1$. Impose exact transport of both potential outcomes. The target covariate mass is either
$p_+=1/2+\varepsilon$ or $p_-=1/2-\varepsilon$.
The ratios relative to the source are bounded, but their values are not revealed. Under $p_+$ selecting $\pi_1$ violates the budget, so its selection probability is at most $\zeta$.

The source observations have identical laws under the two alternatives. Couple each target Bernoulli observation using a common uniform random variable. The probability that any of their $n_0$ coordinates differ is at most $2n_0\varepsilon$. Hence the total variation distance of the available-data laws is at most $2n_0\varepsilon$. Selection of $\pi_1$ under $p_-$ therefore has probability at most $\zeta+2n_0\varepsilon$. But $\pi_1$ is feasible and improves value by $p_-$ there. Taking $\varepsilon=n_0^{-2}$ gives expected regret at least
\[
 (1/2-n_0^{-2})(1-\zeta-2/n_0),
\]
which tends to $(1-\zeta)/2>0$. Learner randomization can be coupled identically and does not change the argument.
\end{proof}

This lower bound deliberately does not give the learner the unknown density ratio: supplying it in this example would reveal $p_+$ or $p_-$. It addresses the target-sampling part of the catalogue problem, while Theorem~\ref{policythm} supplies an explicit achievable benchmark in a more informative experiment.

\section{Extension to learned weights and remaining gap}
There is a direct deterministic robustness statement. If, conditional on independent training, an estimated score $\widetilde Z(\pi)$ is uniformly bounded by a known envelope and satisfies
\[
 \sup_{\pi\in\Pi}|\E_P\widetilde Z(\pi)-v(\pi)|\leq b_n,
\]
then adding $b_n$ to $e_n$ makes the proof of Theorem~\ref{policythm} unchanged. To check this, center the empirical mean at its actual expectation and use the displayed bound before the two value inequalities. This is a certificate requirement, not an assertion that small prediction error automatically supplies it. A substantive next step is to derive such a certificate for estimated transport and treatment weights with outcome augmentation, uniformly over a growing policy class. Sharp dependence on infinite-class complexity, optimal nuisance thresholds, and feasible budget relaxation are not established here.

\begin{thebibliography}{9}
\bibitem{dml} V. Chernozhukov, D. Chetverikov, M. Demirer, E. Duflo, C. Hansen, W. Newey and J. Robins (2018), \emph{Double/debiased machine learning for treatment and structural parameters}, Econometrics Journal 21, C1--C68. Section 1 and the partially linear example in the author-posted version were checked; this paper is background, not a transport theorem used above. \url{https://arxiv.org/pdf/1608.00060}. \url{https://doi.org/10.1111/ectj.12097}.
\bibitem{policy} S. Athey and S. Wager (2021), \emph{Policy Learning with Observational Data}, Econometrica 89, 133--161. Author-posted abstract and bibliographic record inspected 9 October 2026; no specific theorem is invoked. \url{https://arxiv.org/abs/1702.02896}.
\end{thebibliography}

\end{document}
